% Options for packages loaded elsewhere
\PassOptionsToPackage{unicode}{hyperref}
\PassOptionsToPackage{hyphens}{url}
%
\documentclass[
]{article}
\usepackage{amsmath,amssymb}
\usepackage{lmodern}
\usepackage{iftex}
\ifPDFTeX
  \usepackage[T1]{fontenc}
  \usepackage[utf8]{inputenc}
  \usepackage{textcomp} % provide euro and other symbols
\else % if luatex or xetex
  \usepackage{unicode-math}
  \defaultfontfeatures{Scale=MatchLowercase}
  \defaultfontfeatures[\rmfamily]{Ligatures=TeX,Scale=1}
\fi
% Use upquote if available, for straight quotes in verbatim environments
\IfFileExists{upquote.sty}{\usepackage{upquote}}{}
\IfFileExists{microtype.sty}{% use microtype if available
  \usepackage[]{microtype}
  \UseMicrotypeSet[protrusion]{basicmath} % disable protrusion for tt fonts
}{}
\makeatletter
\@ifundefined{KOMAClassName}{% if non-KOMA class
  \IfFileExists{parskip.sty}{%
    \usepackage{parskip}
  }{% else
    \setlength{\parindent}{0pt}
    \setlength{\parskip}{6pt plus 2pt minus 1pt}}
}{% if KOMA class
  \KOMAoptions{parskip=half}}
\makeatother
\usepackage{xcolor}
\setlength{\emergencystretch}{3em} % prevent overfull lines
\providecommand{\tightlist}{%
  \setlength{\itemsep}{0pt}\setlength{\parskip}{0pt}}
\setcounter{secnumdepth}{-\maxdimen} % remove section numbering
\ifLuaTeX
  \usepackage{selnolig}  % disable illegal ligatures
\fi
\IfFileExists{bookmark.sty}{\usepackage{bookmark}}{\usepackage{hyperref}}
\IfFileExists{xurl.sty}{\usepackage{xurl}}{} % add URL line breaks if available
\urlstyle{same} % disable monospaced font for URLs
\hypersetup{
  hidelinks,
  pdfcreator={LaTeX via pandoc}}

\author{Dietmar Gerald Schrausser}
\date{09.05.2026}


\begin{document}

\hypertarget{foliumblat-iiv}{%
\section{Foliu{[}m{]}/Blat IIv}\label{foliumblat-iiv}}

A \emph{contemporary} English translation is presented from comparing
the existing Latin, Early New High German (ENHG) and English
(translation) texts for better comprehensibility. Especially since the
ENHG text is difficult to understand compared to modern High German, as
a \emph{profound} knowledge of German (as a mother tongue) as well as
already aknowledge of various \emph{dialects} is required.

\hypertarget{of-the-work-of-the-first-day}{%
\subsection{Of the Work of the First
Day}\label{of-the-work-of-the-first-day}}

\begin{quote}
\textbf{D}Ixit deus. fiat lux.\\
Et facta e{[}st{]} lux. Et vidit deus lucem q{[}uia{]} esset bona.\\
Et diuisit luce{[}m{]} a tenebris.\\
Appelauitq{[}ue{]} luce{[}m{]} die{[}m{]}.\\
Et tenebras nocte{[}m{]}.\\
Factu{[}m{]}q{[}ue{]} est vespere {[}et{]} mane dies vnus.
\end{quote}

With divine, superhuman diligence, \textbf{Moses} composed a wonderful
masterpiece about all the secrets of nature, a book that surpasses all
others in eloquence and ingenuity\footnote{`eloq{[}ue{]}ntia{[}m{]}
  {[}et{]} ingeniu{[}m{]}' and `außprechlichkeit vn{[}d{]}
  sinnreichikeit'.}: For the glorious God, who \emph{is} the \emph{true
light}, appreciating the light, making all things in the light, most
correctly\footnote{`rectissime' and `gar recht'.} began the creation of
the \emph{world} from the light\footnote{`a luce' and `am liecht'.}.\footnote{attributed
  to Emperor Charlemagne; his reform program `Admonitio Generalis' (s.
  Baluzius,
  \href{https://digital.ub.uni-paderborn.de/eab/content/zoom/1075337}{1677})
  refers exclusively to the Bible and canon law and reflects the
  foundations of the \emph{Carolingian Renaissance}.} This completed a
natural day in its cycle\footnote{`circuito' and `umkreis'.} over three
days until the fourth (in which the luminary\footnote{`lu{[}m{]}inaria
  formata' and `großen liecht geformt'.} are formed).\footnote{from the
  `Etymologiae' (De Seville,
  \href{https://books.google.com/books?id=wb8Z_vlSyJwC}{1802}, Book V,
  Chapter 30).}

Among all physical things, the noblest and, in terms of its spiritual
essence, the closest and optimal, encompassing\footnote{`maxi{[}m{]}e
  co{[}mun{]}icans' and `allermeist gemainsam'.} all in its
beauty\footnote{`pulcritudinem' and `das sein schone'.}, filling even
the smallest point in the world, is exclusively \emph{Light}, which
passes unharmed through impurity, making the whole world good and
beautiful. Deservedly, he saw that the Light which was good was nothing
other than a reflection, a subtle and shadowy likeness of the supreme
Good.\footnote{Passage from the `Heptaplus' (Mirandola \& Mirandola,
  \href{https://doi.org/10.3931/e-rara-61217}{1601}, Book I, Chapter 4).}

As soon as the spirit reached these waters and penetrated the
substructure\footnote{`s{[}u{]}biectu{[}m{]}' and `underwurff'.}, at the
architect's\footnote{`artificis' and `werckmeisters'.} command, light
arose in beauty and splendor, shining like clouds of light illuminating
the upper regions of the world with its radiance (like the approaching
morning sun), first the upper, then subsequently the lower part of the
\emph{sphere}.\footnote{from the `Heptaplus' (Mirandola \& Mirandola,
  \href{https://doi.org/10.3931/e-rara-61217}{1601}, Exposition 1,
  Chapter 2).} Afterward, he separated them, so that darkness and light
would have two distinct hemispheres. Because of its clarity, which
purifies\footnote{`purgat' and `rainigt'.} the darkness, he called the
light day\footnote{`Appellauit luce{[}m{]} die{[}m{]}' c.f. `dimm' and
  `liecht hies er den tag'.} and the harmful darkness night\footnote{`tenebras
  a noce{[}n{]}do nocte{[}m{]}' and `finsternus vo{[}n{]} beschedigung'.},
so that the eyes would not see it\footnote{From a purely physical
  perspective, this seems illogical, especially since excessively bright
  light would likely have harmful effects, which would hardly be the
  case with excessive darkness. Figuratively speaking, however, it seems
  plausible, since \emph{light} in the sense of \emph{knowledge} can
  indeed \emph{illuminate} and thus \emph{purify} \emph{darkness} in the
  sense of \emph{ignorance}.}.\footnote{from the `Historia Scholastica'
  (Comestor,
  \href{https://books.google.com/books?id=fKkLGynZVoYC}{1699}, Liber
  Genesis, Chapter 3).} From the dimensions of these parts, he created
day and night, with the alternating sequence of spaces and cycles
forming the eternal, recurring periods that we call years.\footnote{from
  the `Divinae Institutiones' (Lactantius,
  \href{https://books.google.com/books?id=qGsnghsi3uUC}{1497}, Book II,
  Chapter 9 or 10).}

Thus, \emph{one} day has passed, the first day of the
\emph{world}\footnote{`seculi' (!).}, but not the first of all
days\footnote{`p{[}ri{]}m{[}us{]} dio{[}rum{]} o{[}mn{]}i{[}u{]}m' and
  `der erst aller tag'.}. Therefore, it is not called the first, but
\emph{one} day. Truly\footnote{`sic'.}, on this day God created
formless\footnote{`informe{[}m{]}'.} matter, the angels, the heavens,
the light, the earth, the water, the air (etc.). Furthermore, two
opposite, distinct parts of the earth, the east and the west, were
created.\footnote{from the `Divinae Institutiones' (Lactantius,
  \href{https://books.google.com/books?id=qGsnghsi3uUC}{1497}, Book II,
  Chapter 9 or 10).}

The \emph{East} is attributed to God; the \emph{source}\footnote{`fons'
  and `prun{[}n{]}'.} of light and the illustrator of things, who raises
us to eternal life. The \emph{West}, however, is attributed to the
quarrelsome, indignant and malicious mind\footnote{`{[}con{]}turbate
  illi p{[}ra{]}ueq{[}ue{]} me{[}n{]}ti' and `zerstreitten entrüsten vnd
  boßhafftigen gemüt'.}; it brings darkness, which aims to seduce and
kill people through sin. Just as light springs from the East and worldly
reason\footnote{`vite ratio' and `vernunfft des lebens'.} turns within
the light\footnote{`v{[}er{]}sat{[}ur{]}' and `swebt' (!).}, so darkness
comes from the West; there, death and destruction reside.\footnote{from
  the `Divinae Institutiones' (Lactantius,
  \href{https://books.google.com/books?id=qGsnghsi3uUC}{1497}, Book II,
  Chapter 9 or 10).}

Then God determined the other regions, the south and the north, in the
same dimension, corresponding to the first two parts. Thus, the region
strongly warmed by the sun is closest to the east. Conversely, that
which is permanently frozen and numb with cold corresponds to a part of
the extreme west\footnote{`cui{[}us{]} extrem{[}us{]} occasus' and `des
  letzte{[}n{]} niderga{[}n{]}gs'.}. As darkness is opposite to light,
so cold is opposite to warmth. Just as warmth is closest to light, so
the south is closest to the east and just as cold is closest to
darkness, so too is the plague of the north\footnote{`plaga
  septe{[}m{]}trionis'.} closest to the west. All this is clarified in
the \emph{Work of the Fourth Day}.\footnote{from the `Divinae
  Institutiones' (Lactantius,
  \href{https://books.google.com/books?id=qGsnghsi3uUC}{1497}, Book II,
  Chapter 9 or 10).}

\hypertarget{references}{%
\subsection{References}\label{references}}

Baluzius, S., ed.~(1677). Capitvlaria Karoli Magni. In \emph{Capitvlaria
Regvm Francorvm}, 1:189. Parisiis: Muguet.
\url{https://digital.ub.uni-paderborn.de/eab/content/zoom/1075337}

Comestor, P. (1699). \emph{Eruditissimi Viri Magistri Petri Comestoris
Historia Scholastica: Opus Eximium Magnam Sacrae Scripturae Partem, Quae
{[}Et{]} in Serie, {[}Et{]} in Glossis Crebro Diffusa Erat, Breviter
Complectens}. Matriti: ex officina Antonij Gonçalez de Reyes \ldots{} :
a costa de Francisco Sacédon \ldots{}
\url{https://books.google.com/books?id=fKkLGynZVoYC}

De Seville, I. (1802). \emph{S. Isidori Hispalensis Episcopi Hispaniarum
Doctoris Opera Omnia}. Edited by Arevalo, F. Romae: Apvd Antonivm
Fvlgonivm. \url{https://books.google.com/books?id=wb8Z_vlSyJwC}

Lactantius, L.C.F. (1497). \emph{Divinae Institutiones}. Venetiis: Simon
Bevilaqua. \url{https://books.google.com/books?id=qGsnghsi3uUC}

Mirandola, G., \& Mirandola, G. F. (1601). \emph{Ioannis Pici,
Mirandulae \ldots{} opera quae extant omnia \ldots{}} Basileae: per
Sebastianum Henricpetri. \url{https://doi.org/10.3931/e-rara-61217}

\end{document}
